\section{Projective residue states and static field separation}
\label{sec:projective}

\begin{definition}[Projective residue block]
For $n\ge2$, a pair $(a,b)\in\Rring_n^2$ is unimodular if there exist
$u,v\in\Rring_n$ such that $ua+vb=1$.  Define
\[
 X_n=P^1(\Rring_n)
 =\{(a,b):\Rring_na+\Rring_nb=\Rring_n\}/\Rring_n^\times,
\]
where a unit acts by simultaneous multiplication of both coordinates.
\end{definition}

This definition uses no field test.  In particular, it is valid for prime
powers and mixed composite moduli, where zero divisors alter the geometry.
The unit-scaling quotient is also compatible with a transported semiring
presentation, a property used in \cref{sec:exact-audit}.

\begin{lemma}[Prime-power count]
\label{lem:prime-power-count}
For $q=p^a$,
\[
 |P^1(\Z/q\Z)|=q+q/p.
\]
\end{lemma}

\begin{proof}
The ring $A=\Z/q\Z$ is local with maximal ideal $pA$.  A unimodular pair has
at least one unit coordinate.  If the second coordinate is a unit, unit
scaling gives a unique representative $[t:1]$ with $t\in A$, hence $q$
classes.  If the second coordinate is a nonunit, the first is a unit and the
class has a unique representative $[1:u]$ with $u\in pA$, hence $q/p$
further classes.  The two families are disjoint.
\end{proof}

\begin{lemma}[Product compatibility]
\label{lem:product-projective}
For finite commutative rings $A,B$,
\[
 P^1(A\times B)\cong P^1(A)\times P^1(B).
\]
\end{lemma}

\begin{proof}
Unimodularity and unit scaling are componentwise over $A\times B$.
Passing to equivalence classes gives the bijection.
\end{proof}

\begin{theorem}[Projective count]
\label{thm:projective-count}
For every $n\ge2$,
\[
 |X_n|=\psi(n):=n\prod_{p\mid n}\left(1+\frac1p\right).
\]
\end{theorem}

\begin{proof}
Apply the Chinese remainder theorem to the prime-power factorization of $n$
and combine \cref{lem:prime-power-count,lem:product-projective}.
\end{proof}

\begin{corollary}[Static field criterion]
\label{cor:static-field}
For $n\ge2$,
\[
 |X_n|=n+1\quad\Longleftrightarrow\quad n\text{ is prime}.
\]
\end{corollary}

\begin{proof}
For $n=p$, \cref{thm:projective-count} gives $p+1$.  Conversely, if $n$ is
composite and $p\mid n$, then $n/p\ge2$ and
\[
 \psi(n)\ge n(1+1/p)=n+n/p\ge n+2.
\]
\end{proof}

The defect
\[
 \delta(n):=|X_n|-(n+1)
\]
therefore recognizes precisely when the residue ring is a field.  This is a
real source-level arithmetic distinction, and it escapes the bare
multiplicative-UFD clone inherited from the earlier source program.  It is
not yet a primitive-orbit distinction.  Every state and transition in the
next section is created without consulting $\delta(n)$; only that frozen
ledger can decide A1.
